\subsection{An actual stack realization of the normalized word trace}
The modern stack formalism changes the finiteness premise of
\fid{F26}. The primary inputs used here are Sun's trace formula
\cite[Theorem~1.1, p.~49; Definition~4.1, p.~84; Theorem~4.2,
p.~85]{SunStacks}, his explicit classifying-stack calculation
\cite[Theorem~4.8 and \S4.8.2, pp.~89--90; Example~7.4,
p.~109]{SunStacks}, and the published Laszlo--Olsson duality and
K\"unneth formulas \cite[\S9.1, Lemmas~9.1.1--9.1.2, p.~198;
Theorem~11.0.14, p.~204]{LO2}. These statements and their applicable
hypotheses were inspected; their full general proofs are cited inputs.

\begin{finding}[\fid{F35}: a convergent stack trace with an unbounded dual]\label{f:stacktrace}
Over $\F_q$ let
\[
 \mathcal Y=[(\A^6\times\mathbb G_m^8)/(\GL_2)^4]
   =\A^6\times\mathbb G_m^8\times(B\GL_2)^4,
\]
with trivial action, and let $\mathcal X$ be the disjoint union
of 24 rational points and 16 copies of $\mathcal Y$.
Then $\mathcal X$ is a finite-type Artin stack with affine stabilizers,
and for every $n\ge1$, with $Q=q^n$,
\begin{equation}\label{eq:stackwordtrace}
 \#\mathcal X(\F_Q)
 =24+\frac{16Q^2}{(Q+1)^4}
 =\sum_i(-1)^i\Tr(F^n\mid H_c^i(\mathcal X_{\bar{\F}_q},\Q_\ell)).
\end{equation}
Here $\#$ means groupoid mass and the series is absolutely convergent.
This realizes the complete scalar sequence of \fid{F26}.

The component $\mathcal Y$ is smooth of dimension $-2$.
Its compact cohomology is unbounded below and satisfies actual duality
\[
 H_c^{-4-i}(\mathcal Y)\simeq H^i(\mathcal Y)^\vee(2).
\]
Its visible compact spectrum is bounded and admissible as in
\fid{F32}; the corresponding ordinary dual spectrum is unbounded.
Thus actual duality holds, while the separate dual upper-bound
hypothesis of \fid{F17}/\fid{F32} fails in this model.
No nonstandard generator action or simultaneous word comparison
is asserted.
\end{finding}

\begin{proof}
Put $U=\A^6\times\mathbb G_m^8$ and $G=(\GL_2)^4$.
The atlas $U\to[U/G]$ is smooth of relative dimension 16, whereas
$\dim U=14$. This gives the stated dimension and smoothness of
$\mathcal Y$. The full $\mathcal X$ is not equidimensional: its
components have dimensions zero and minus two, and $\dim\mathcal X=0$.

A $\GL_2$-torsor over a field is the frame torsor of a rank-two
vector space, hence is trivial. For the trivial action each
$U(\F_Q)$ point therefore contributes $|G(\F_Q)|^{-1}$ to the mass.
Counting linearly independent columns gives
\[
 |\GL_2(\F_Q)|=(Q^2-1)(Q^2-Q)=Q(Q-1)^2(Q+1).
\]
Consequently
\[
 \#\mathcal Y(\F_Q)
 =\frac{Q^6(Q-1)^8}{Q^4(Q-1)^8(Q+1)^4}
 =\frac{Q^2}{(Q+1)^4}.
\]
Sun's Theorem~1.1 gives the asserted convergent trace for the constant
sheaf. Equivalently Theorem~4.2 applies: the constant sheaf is
bounded, geometrically stratifiable, pure of weight zero, and
convergent in Definition~4.1. The structural map is finite type.
The finite-field base also satisfies the excellence, invertibility
of $\ell$, and finite cohomological-dimension hypotheses of the
cited Laszlo--Olsson formalism.

We now compute the actual cohomology, suppressing the geometric
base change and coefficient field in the notation. Define
\[
 \sum_{s\ge0}a_s z^s=(1-z)^{-4}(1-z^2)^{-4}.
\]
For one classifying factor the known polynomial presentation is
$H^*(B\GL_2)=\Q_\ell[c_1,c_2]$, with $\deg c_i=2i$ and
$c_i$ a rational Chern class spanning $\Q_\ell(-i)$.
Sun's Example~7.4 explicitly gives
\[
 H_c^{-8-2s}(B\GL_2)
 =\Q_\ell(s+4)^{\oplus(\lfloor s/2\rfloor+1)},\qquad s\ge0.
\]
The ordinary cohomology of $\mathbb G_m^8$ is the exterior
algebra on eight copies of $\Q_\ell(-1)$ in degree one, while
$\A^6$ has only the unit in ordinary cohomology. K\"unneth gives
\begin{equation}\label{eq:stackordinarygroups}
 H^i(\mathcal Y)=
 \bigoplus_{\substack{0\le b\le8,\ s\ge0\\i=b+2s}}
 \Q_\ell(-(b+s))^{\oplus\binom8b a_s}.
\end{equation}
For a direct compact-support calculation, the affine factor has
$H_c^{12}(\A^6)=\Q_\ell(-6)$, the torus has
$H_c^{16-b}(\mathbb G_m^8)=\Q_\ell(-(8-b))^{\oplus\binom8b}$,
and the four classifying factors give
$H_c^{-32-2s}((B\GL_2)^4)=\Q_\ell(s+16)^{\oplus a_s}$.
The compact-support K\"unneth formula therefore yields
\begin{equation}\label{eq:stackcompactgroups}
 H_c^j(\mathcal Y)=
 \bigoplus_{\substack{0\le b\le8,\ s\ge0\\j=-4-b-2s}}
 \Q_\ell(b+s+2)^{\oplus\binom8b a_s}.
\end{equation}
The sum in each degree is finite. Since $a_s>0$ for every $s$,
these groups are unbounded below. The indicated weight is
$-2(b+s+2)=j-b$, consistent with Sun's affine-stabilizer bound
\cite[Theorem~1.4, p.~50]{SunStacks}. A fixed degree can be mixed:
$H_c^{-6}=\Q_\ell(3)^4\oplus\Q_\ell(4)^{28}$.

For $x=q^{-n}$ the sign in \eqref{eq:stackcompactgroups} is
$(-1)^b$, and the Euler series becomes
\[
 x^2\frac{(1-x)^8}{(1-x)^4(1-x^2)^4}
 =\frac{x^2}{(1+x)^4}.
\]
This grouping is legitimate: for every real $t>0$ the ungrouped
absolute eigenvalue moment is
\[
 \sum_{b,s}\binom8b a_s q^{-t(b+s+2)}
 =\frac{q^{-2t}(1+q^{-t})^4}{(1-q^{-t})^8}<\infty.
\]
The visible coefficient of $q^{-(j+2)}$ is consequently
\[
 \sum_{b=0}^{\min(8,j)}(-1)^b\binom8b a_{j-b}
   =(-1)^j\binom{j+3}{3}.
\]
For instance at $q^{-3}$ the actual degree-$-6$ multiplicity four
and degree-$-5$ multiplicity eight give the signed coefficient
$4-8=-4$. The full $\mathcal X$ therefore has exactly the
admissible visible spectrum
\[
 (1,24),\qquad
 \left(q^{-(j+2)},16(-1)^j\binom{j+3}{3}\right),\quad j\ge0.
\]
Its support and absolute moments meet \fid{F32}'s hypotheses, so
that finding's uniqueness applies in its stated convergence class.

Finally the smooth dualizing complex of $\mathcal Y$ is
$\Q_\ell(-2)[-4]$. The cited duality formula gives
$H_c^{-4-i}=H^i{}^\vee(2)$, also verified term by term by
\eqref{eq:stackordinarygroups}--\eqref{eq:stackcompactgroups}.
A compact root $\alpha=q^{-(b+s+2)}$ is paired with the
ordinary root $q^{-2}/\alpha=q^{b+s}$. After Euler grouping
the latter roots are all $q^j$, with nonzero coefficients
$(-1)^j\binom{j+3}{3}$. Their support is unbounded, and the
ordinary geometric-Frobenius trace terms do not tend to zero.
Thus this dual is outside the independently bounded visible-spectrum
class. For the nonequidimensional $\mathcal X$, Verdier duality
is computed componentwise; no common dimension-minus-two pairing
is imposed on all its components.
\end{proof}

This is a uniform geometric realization of one scalar word-trace
sequence, with unbounded cohomology explaining its compatibility
with \fid{F26}. It supplies neither operators for $e_1,e_2$ nor
their composition law, and no tensor-path sign or nonstandard basis
follows from its groupoid mass. Its failure of the independent dual
bound is a proved property of this model, not an exclusion of every
stack or completed realization.

% Separate paragraph for Part III.
\paragraph{\fid{F35} in Grothendieck's language.}
The ordinary split group schemes and quotient presentation are
explicit. Positive-dimensional stabilizers are retained, so this
construction uses the Artin-stack extension of the six operations,
not the bounded constructible pushforward theorem for schemes.
Its rational mass and infinite alternating Tate trace are connected
by the precisely cited Laszlo--Olsson/Sun formalism. The effective
cohomology groups in \eqref{eq:stackcompactgroups} have positive
dimensions, while their grouped Euler coefficients have both signs.
The duality exists but does not preserve the chosen convergent
geometric-Frobenius trace class. The audited manuscript determinant
and projector vocabulary is only a framework anchor; no specific
stack, nonstandard action, or proof here is attributed to those leaves.

% New bibliography keys:
% SunStacks: S. Sun, L-series of Artin stacks over finite fields,
% Algebra & Number Theory 6 (2012), 47--122.
% https://msp.org/ant/2012/6-1/ant-v6-n1-p03-s.pdf
% LO2: Y. Laszlo and M. Olsson, The six operations for sheaves on Artin
% stacks II: Adic coefficients, Publ. Math. IHES107(2008),169--210.
% https://numdam.org/item/10.1007/s10240-008-0012-5.pdf
